\section{Tokenizer}
\label{sec:tokenizer}

The tokenizer is a Byte-Pair Encoding (BPE) model trained on a combined corpus
using the HuggingFace \texttt{tokenizers} library.  The vocabulary size is
\textbf{50{,}000} sub-word tokens, stored in
\texttt{tokenizers/tokenizer\_corpus.json}.

\textbf{Design choices:}
\begin{itemize}[nosep]
  \item \textbf{BPE over WordPiece or Unigram:}  BPE provides a good balance
        between vocabulary efficiency and handling of out-of-vocabulary tokens.
        It is also the standard for GPT-family models.
  \item \textbf{Corpus-wide training:}  The tokenizer is trained on a combined
        corpus spanning all domains (TinyStories, SimpleWiki, WritingPrompts,
        etc.), ensuring stable tokenization across stages.  A tokenizer trained
        only on children's stories would produce poor segmentation on adult
        fiction vocabulary.
  \item \textbf{Special tokens:}  \texttt{BOS} (id=2) and \texttt{EOS} (id=3)
        delimit documents during chunking.
  \item \textbf{Frozen across stages:}  The same tokenizer is used for all
        six stages.  Changing the tokenizer mid-training would invalidate the
        embedding matrix and cached datasets.
\end{itemize}

\textbf{Vocabulary size trade-off.}  50K tokens is a compromise: large enough
to cover the combined vocabulary with reasonable sub-word granularity, but
small enough that the embedding layer ($50000 \times 512 = 25.6$M parameters
with weight tying) does not dominate the parameter budget.
